\documentclass[11pt,a4paper]{article}
\usepackage{turkos}

\begin{document}
\phaseheader{5}{Physical Memory Management}{Find out which RAM really exists, and hand it out and take it back in 4 KiB page frames}{28}{34}

\ataglance{1--2 weeks}{3}{Phase 4: working interrupts, timer, keyboard monitor}{A parser for the Multiboot memory map,
linker symbols that mark the kernel's start and end, a bitmap page-frame allocator with \code{pmm_alloc_frame},
\code{pmm_free_frame} and contiguous allocation, a \code{mem} command with real numbers, and the first in-kernel
test suite.}

\section{Why this phase}

Every later subsystem needs memory: page tables (Phase 6), the kernel heap (Phase 7), process stacks (Phase 8), user
address spaces (Phase 10), the buffer cache (Phase 12). All of it ultimately comes from one place, the
\textbf{physical memory manager} (PMM), which knows which physical page frames are free and hands them out one at a
time.

The PMM has to answer a harder question first: which memory can it hand out at all? Physical addresses are not a
clean array of RAM. The first megabyte is full of BIOS data and video memory, the kernel image sits at 1\,MiB,
GRUB's structures and boot modules are scattered above it, and there are holes near the top of the 32-bit space for
devices and firmware tables. Hand out the wrong frame and you will overwrite the kernel, or a device, and the crash
will come much later somewhere unrelated. This phase is about getting that bookkeeping exactly right. It is also
where you start writing \textbf{automated tests}, because a memory allocator that is ``mostly right'' is wrong.

\section{Concepts you need}

\subsection{Address spaces and why the kernel needs an allocator}

\OSTEP chapter 13 introduces the \textbf{address space}, the memory a program believes it has, and explains why the
OS must virtualise memory: isolation, protection, and the convenience of every program seeing memory starting at
zero. You will build that illusion in Phases 6 and 10. Underneath it, the OS still manages real memory in fixed-size
units called \textbf{page frames}. On x86 a frame is 4\,KiB, and frame number $n$ covers physical addresses
$n \times 4096$ to $n \times 4096 + 4095$. Frame number and address convert with a shift by 12: \code{frame = addr >> 12}.

\subsection{The Multiboot memory map}

Because you set bit 1 of the Multiboot flags in Phase 2, GRUB passes a \textbf{memory map} (\code{mmap_addr},
\code{mmap_length} in the Multiboot information, valid when bit 6 of \code{flags} is set). It is the BIOS E820 map in
a friendlier wrapper: a list of regions, each with a 64-bit base address, a 64-bit length and a type. Only type 1 is
usable RAM; everything else (reserved, ACPI tables, bad memory) must never be allocated. Figure~\ref{fig:physmap}
shows a typical layout for a 128\,MiB QEMU machine and what the PMM does with it.

\begin{figure}[htbp]
\centering
\begin{tikzpicture}[m/.style={draw=tkNavy!70, minimum width=6.8cm, font=\sffamily\scriptsize, align=center},
  a/.style={font=\ttfamily\scriptsize, anchor=east}, r/.style={font=\sffamily\scriptsize, anchor=west, text=tkGray}]
  \node[m, minimum height=0.6cm, fill=tkGray!18] (top) at (0,0) {reserved: firmware, device memory, ACPI (never allocated)};
  \node[m, minimum height=2.2cm, fill=tkGreen!12, below=0pt of top] (free) {\textbf{free frames}\\handed out by \code{pmm_alloc_frame()}};
  \node[m, minimum height=0.55cm, fill=tkAmber!15, below=0pt of free] (mods) {Multiboot info, memory map, boot modules};
  \node[m, minimum height=0.55cm, fill=tkBlue!12, below=0pt of mods] (bm) {PMM bitmap (placed right after the kernel)};
  \node[m, minimum height=0.9cm, fill=tkRed!10, below=0pt of bm] (k) {\textbf{kernel image}: \code{.text .rodata .data .bss}};
  \node[m, minimum height=0.7cm, fill=tkGray!18, below=0pt of k] (low) {first MiB: BIOS, VGA, ROMs (see Phase 1)};
  \node[a] at ([xshift=-4pt]top.north west) {0x08000000};
  \node[a] at ([xshift=-4pt]top.south west) {\textasciitilde 0x07FE0000};
  \node[a] at ([xshift=-4pt]bm.north west) {};
  \node[a] at ([xshift=-4pt]k.north west) {\code{kernel_end}};
  \node[a] at ([xshift=-4pt]k.south west) {0x00100000};
  \node[a] at ([xshift=-4pt]low.south west) {0x00000000};
  \node[r] at ([xshift=4pt]top.east) {type 2 in the memory map};
  \node[r] at ([xshift=4pt]free.east) {type 1};
  \node[r, align=left] at ([xshift=4pt]mods.south east) {type 1, but in use:\\mark as used};
  \node[r] at ([xshift=4pt]low.east) {mixed; never allocate};
\end{tikzpicture}
\caption{Physical memory of a 128\,MiB QEMU machine (not to scale). Only frames in the green area are free.}
\label{fig:physmap}
\end{figure}

\subsection{Tracking free frames: bitmap or free list}

There are two classic ways to track which frames are free, compared in \MOS\,\S3.2.3, \OSDI\,\S4.2 and \OSTEP
chapter 17. A \textbf{bitmap} uses one bit per frame. It is compact, and it is easy to find a run of free frames for a
contiguous allocation, but finding a free frame means scanning. A \textbf{free list} chains free frames together,
often by storing the ``next'' pointer inside each free frame itself, which gives O(1) allocation but makes contiguous
runs hard to find. The table puts numbers on the choice.

\begin{center}
\small\renewcommand{\arraystretch}{1.3}
\begin{tabularx}{\linewidth}{@{}l X X@{}}
\toprule
& \tkth{Bitmap} & \tkth{Free list (stack of frames)} \\
\midrule
Memory overhead & 1 bit per 4\,KiB frame: 128\,KiB for all 4\,GiB & none extra: the list lives in the free frames \\
Allocate one frame & scan for a 0 bit; fast with a ``next free'' hint & pop the head: O(1) \\
Free one frame & clear one bit: O(1) & push on the head: O(1) \\
Contiguous run of $n$ frames & scan for $n$ zero bits in a row & impractical \\
Detect double free & trivial: the bit is already 0 & needs extra bookkeeping \\
\bottomrule
\end{tabularx}
\end{center}

Turk-OS starts with a bitmap: it makes bugs visible, supports the contiguous allocations that DMA needs in Phase 12,
and reuses your Phase 0 bitmap code. You can add a free-list cache in front of it later if profiling says so.

\section{Reading plan}

\begin{readingplan}
\must & \LOB & Ch.~10 \emph{Page Frame Allocation}: how much memory is there, managing available memory & 61--63 & Tasks 5.1--5.3. \\
\must & \OSTEP & Ch.~13 \emph{The Abstraction: Address Spaces} & 113--122 & The goals that all of memory management serves. \\
\must & \OSTEP & Ch.~17 \emph{Free-Space Management} & 161--178 & Bitmaps, free lists, splitting, coalescing; you will reuse this in Phase 7. \\
\should & \OSTEP & Ch.~14 \emph{Interlude: Memory API} & 123--134 & The bugs allocator users make, so your tests look for them. \\
\should & \MOS & \S3.1 \emph{No Memory Abstraction}; \S3.2 \emph{Address Spaces} (esp.\ \S3.2.3 \emph{Managing Free Memory}) & 182--194 & Bitmaps versus linked lists with worked examples. \\
\should & \OSDI & \S4.1--4.2 \emph{Basic Memory Management}, \emph{Swapping} (bitmaps, linked lists) & 374--383 & The same ideas, with the MINIX perspective. \\
\could & \OSC & \S9.1--9.2 \emph{Background}; \emph{Contiguous Memory Allocation} & 349--360 & Fragmentation, first/best/worst fit. \\
\could & \OSDI & \S4.7.1 \emph{Memory Layout} (MINIX process manager) & 422--425 & How MINIX lays out memory for processes. \\
\end{readingplan}

\begin{minixbox}
MINIX 3 as printed in \OSDI does not use paging at all: the process manager allocates each process a contiguous
chunk of physical memory from a list of \emph{holes}. The hole list is explained on pp.~431 and 473--474 of \OSDI;
the file that implements it, \texttt{servers/pm/alloc.c}, is described there but not included in the printed listing. That is the linked-list scheme from \MOS\,\S3.2.3 used for whole processes. Turk-OS uses
paging instead, so it only ever needs single 4\,KiB frames, apart from the occasional contiguous DMA buffer. The
comparison shows why paging won: no external fragmentation, and no need to move processes around.
\end{minixbox}

\section{Build it}

\task{Walk the memory map}
Extend \code{multiboot.h} with the memory map entry and print the map at boot. The one trap: each entry starts with a
\code{size} field that does \emph{not} count itself, so the next entry is \code{size + 4} bytes further on.

\begin{lst}{c}{kernel/mm/pmm.c (memory map walk)}
struct mmap_entry {
    uint32_t size;          /* size of the rest of this entry, not counting this field */
    uint64_t base;
    uint64_t length;
    uint32_t type;          /* 1 = usable RAM */
} __attribute__((packed));

void print_memory_map(const struct multiboot_info *mbi)
{
    ASSERT(mbi->flags & (1 << 6));                  /* memory map present */
    uintptr_t p   = mbi->mmap_addr;
    uintptr_t end = mbi->mmap_addr + mbi->mmap_length;
    while (p < end) {
        const struct mmap_entry *e = (const struct mmap_entry *)p;
        kprintf("  %08x%08x  len %08x%08x  type %u\n",
                (uint32_t)(e->base >> 32), (uint32_t)e->base,
                (uint32_t)(e->length >> 32), (uint32_t)e->length, e->type);
        p += e->size + sizeof(e->size);
    }
}
\end{lst}
\verify{With \code{-m 128}, you see a usable region from 0 to about 640\,KiB, reserved regions around the BIOS area,
and a large usable region from 1\,MiB to just under 128\,MiB.}

\task{Know where the kernel ends}
The linker can define symbols at any point in the layout. Add \code{kernel_start} and \code{kernel_end} to
\code{link.ld}; in C, declare them as arrays so you use their \emph{address}, never their contents.

\begin{lst}{plain}{link.ld (changes)}
SECTIONS
{
    . = 0x00100000;
    kernel_start = .;
    /* ... .text .rodata .data .bss as before ... */
    kernel_end = .;
}
\end{lst}

\begin{lst}{c}{kernel/mm/pmm.c}
extern char kernel_start[], kernel_end[];     /* addresses only: never read them */
kprintf("kernel: %p - %p (%u KiB)\n", kernel_start, kernel_end,
        (uint32_t)(kernel_end - kernel_start) / 1024);
\end{lst}

\task{The bitmap allocator}
Write \code{pmm_init(mbi)} in four steps. First, find the highest usable address to size the bitmap, and place the
bitmap right after \code{kernel_end}, page-aligned. Second, mark \emph{every} frame as used. Third, walk the memory
map and free the frames of every type-1 region, rounding the start up and the end down to whole frames. Fourth,
mark as used again everything below 1\,MiB, the kernel image, the bitmap itself, the Multiboot information and memory
map, and every boot module (\code{mods_addr}). Starting from ``all used'' means a mistake leaves memory unused instead
of handing out memory that is in use.

\begin{lst}{c}{include/pmm.h}
#define PAGE_SIZE 4096

void      pmm_init(const struct multiboot_info *mbi);
uintptr_t pmm_alloc_frame(void);            /* returns a physical address, 0 if none free */
uintptr_t pmm_alloc_frames(size_t count);   /* physically contiguous run                  */
void      pmm_free_frame(uintptr_t paddr);
size_t    pmm_free_count(void);
size_t    pmm_total_count(void);
\end{lst}

\begin{lst}{c}{kernel/mm/pmm.c (allocation core)}
static uint32_t *bitmap;          /* bit set = frame used */
static size_t    nframes, nfree, next_hint;

uintptr_t pmm_alloc_frame(void)
{
    for (size_t i = 0; i < nframes; i++) {
        size_t f = (next_hint + i) % nframes;
        if (!bitmap_test(bitmap, f)) {
            bitmap_set(bitmap, f);
            nfree--;
            next_hint = f + 1;
            return (uintptr_t)f * PAGE_SIZE;
        }
    }
    return 0;                     /* out of memory: frame 0 is never free, so 0 is safe */
}

void pmm_free_frame(uintptr_t paddr)
{
    ASSERT((paddr & (PAGE_SIZE - 1)) == 0);            /* must be frame-aligned */
    size_t f = paddr / PAGE_SIZE;
    ASSERT(f < nframes && bitmap_test(bitmap, f));     /* catches double frees  */
    bitmap_clear(bitmap, f);
    nfree++;
}
\end{lst}
\verify{After \code{pmm_init}, print total and free frames. The free count should equal the usable RAM above
the bitmap minus the Multiboot structures, give or take a few frames. \code{pmm_alloc_frame()} never returns an
address inside the kernel.}

\task{A tiny test framework}
Write \code{KTEST(name)} and \code{KASSERT(cond)} macros and a \code{test} monitor command that runs every
registered test and prints \code{PASS} or \code{FAIL} with the failing line. A simple way to register tests is an
array of \code{\{name, function\}} pairs in \code{tests/tests.c}; a neater way is to put a pointer to each test in a
dedicated linker section and walk it from \code{__start_ktests} to \code{__stop_ktests}. Phase 7 connects this to
\code{make test} so it runs without you.

\task{Test the allocator hard}
Write at least these tests. Each one should leave the free count exactly where it found it.

\begin{center}
\small\renewcommand{\arraystretch}{1.3}
\begin{tabularx}{\linewidth}{@{}l X@{}}
\toprule
\tkth{Test} & \tkth{What it checks} \\
\midrule
\code{alloc_free_one} & allocate a frame, check alignment and range, free it, count restored \\
\code{no_duplicates} & allocate 1000 frames into an array, check that no two are equal, free all \\
\code{exhaust} & allocate until 0 is returned, check the count reached zero, free everything, count restored \\
\code{never_kernel} & during \code{exhaust}, check that no frame overlaps \code{[kernel_start, kernel_end)} or the bitmap \\
\code{contiguous} & \code{pmm_alloc_frames(16)} returns 16 adjacent frames; free them one by one \\
\code{random_pattern} & 10\,000 random alloc/free steps against a shadow array; counts must match at every step \\
\bottomrule
\end{tabularx}
\end{center}
\verify{All tests pass with \code{-m 16M}, \code{-m 128M}, \code{-m 512M} and \code{-m 3G}. The last one has a
memory hole below 4\,GiB, so it exercises the memory map handling.}

\task{A real mem command}
Update the monitor's \code{mem} command to print the memory map, the kernel's size, and total, used and free
memory in MiB. You will watch these numbers constantly in later phases to spot leaks.

\section{Testing and debugging}

\begin{pitfalls}
The memory map loop runs forever or prints garbage & Advanced by \code{sizeof(struct mmap_entry)} instead of \code{size + 4}, or struct not packed & Use the \code{size} field; check \code{sizeof} is 24. \\
Wrong totals on machines with more than 4\,GiB & Base and length are 64-bit and were truncated & Keep them \code{uint64_t}; clamp regions to the 4\,GiB you can address. \\
Crash long after boot, in code unrelated to the PMM & A frame inside the kernel, the bitmap, or the Multiboot data was handed out and overwritten & Run the \code{never_kernel} test; re-reserve those regions after freeing type-1 memory. \\
Off-by-one: last frame never allocated, or one frame past the end & Region end rounded the wrong way, or \code{<=} instead of \code{<} & Round the start up and the end down; write a test for a region ending exactly on a frame boundary. \\
Allocation gets slow as memory fills & Scanning from frame 0 every time & Keep a \code{next_hint}; later, scan 32 frames at a time by comparing whole words with \code{0xFFFFFFFF}. \\
\end{pitfalls}

\begin{warnbox}[Prepare for Phase 6]
Right now physical address equals the address your code uses, so the PMM can read and write its bitmap directly.
Once paging is on and the kernel moves to \code{0xC0000000}, that stops being true: the bitmap has a physical address
and a different virtual one, and a frame returned by \code{pmm_alloc_frame} can't be touched until it is mapped.
Keep all address arithmetic in the PMM explicit about which kind of address it is, for example by using a
\code{typedef uintptr_t paddr_t} for physical addresses.
\end{warnbox}

\section{Milestone}

\begin{milestonebox}[The kernel knows its memory]
\begin{checklist}
  \item The memory map is printed at boot, and \code{mem} shows totals that match QEMU's \code{-m}.
  \item \code{kernel_start} and \code{kernel_end} come from the linker script.
  \item \code{pmm_alloc_frame}, \code{pmm_free_frame} and \code{pmm_alloc_frames} work, with asserts for alignment and double frees.
  \item All six PMM tests pass from 16\,MiB to 3\,GiB of RAM.
  \item The \code{test} command runs every registered test and reports the results.
\end{checklist}
\end{milestonebox}

\begin{questionbox}
\begin{enumerate}
  \item Why can't you assume that all memory from 1\,MiB to \code{mem_upper} is free?
  \item How big is a bitmap for 4\,GiB of RAM with 4\,KiB frames? For 128\,MiB?
  \item Why start by marking every frame as used, rather than every frame as free?
  \item Why are \code{kernel_start} and \code{kernel_end} declared as arrays in C?
  \item What is the difference between internal and external fragmentation? Which one can a page-frame allocator suffer from?
  \item What would you need to change to make \code{pmm_alloc_frame} O(1)?
\end{enumerate}
\end{questionbox}

\section{Going further}

Implement a \textbf{buddy allocator} (\OSC\,\S10.8.1, \OSTEP\,\S17.4) for power-of-two runs of frames and compare it
with your bitmap. Split memory into zones, for example frames below 16\,MiB for old DMA devices and the rest for
everything else, as Linux does. Or add a small per-CPU cache of free frames in front of the bitmap, which is the
standard way to make the common case O(1).

\nextphase{6}{Paging and the Higher-Half Kernel}
\booklegend
\end{document}
